% Section 8: Through-crossings, resolution, and monodromy
\label{sec:through-body}

\S\ref{sec:classical} produced local monodromy operators and then observed
(No-go~\ref{nogo:unipotent}) that a single one of them carries no spectrum.  This
section supplies the missing structure.  Its central claim is a change of genre: passing
through a singular fibre is not an arrow but a cospan together with a choice, and once
that is made explicit, monodromy becomes a \emph{difference} of the two choices rather
than a primitive.  The change pays immediately: the Legendre monodromy is derived from
a quadratic base change, and a primitive transvection is realised on a finite cell
complex.

\subsection{Crossing is a cospan with a resolution}

\begin{definition}[Through-crossing]\label{def:through-crossing}
Let $F_0$ be a singular fibre and $F_{\pm}$ the two smooth fibres on either side of the
critical value.  A \emph{through-crossing} is the cospan
\[
  F_-\xrightarrow{\ \pi_- \ } F_0 \xleftarrow{\ \pi_+ \ } F_+ ,
\]
whose \emph{through relation} is the fibre product
\[
  R_c=F_-\times_{F_0}F_+ .
\]
\end{definition}

\begin{proposal}[A crossing is not a transport arrow]\label{prop:crossing-not-arrow}
There is no canonical smooth transport arrow $F_-\to F_+$ across a singular fibre.
What exists is the through relation $R_c$ together with a further choice, and the choice
is the mathematical content.
\end{proposal}

\begin{definition}[Resolution torsor]\label{def:resolution-torsor}
A \emph{resolution} of the through-crossing is a selection, for each point of the
critical fibre, of which smooth point on the negative side is joined to which on the
positive side.  The set of resolutions is a torsor under the monodromy group generated
by the vanishing transvection:
\[
  \mathrm{Res}(R_c)\ \text{is a torsor under}\ \langle T_\delta\rangle\cong\Z .
\]
\end{definition}

\begin{remark}[Why the torsor language is not decoration]
A resolution has no privileged origin: any two resolutions differ by a power of the
transvection, and no resolution is singled out by the geometry of the cospan alone.  This
is the precise sense in which \S\ref{sec:classical}'s monodromy operators are not
intrinsic to a single fibre.
\end{remark}

\begin{remark}[A worked affine instance, imported from Paper 0]
\label{rem:affine-instance-of-the-torsor}
The resolution torsor of Definition~\ref{def:resolution-torsor} is an abstract
statement until an instance is produced.  Paper~0 supplies one in the affine sector,
and we cite it rather than re-derive it.  There,
\(\operatorname{Aff}^{+}(1,\mathbb R)\) acts on \(\mathbb H^2\); a dilation is
hyperbolic and its conjugacy class is determined by its multiplier, so its
translation length is a function of the class; a nontrivial translation is
parabolic, and by the two-layer proposition of Paper~0
\cite{Yuan2026AEGPaperZero} every nontrivial translation of a fixed sign is
conjugate to every other, the stabiliser being the translation subgroup.  The class of a translation therefore carries no
continuous invariant, and the set of its conjugates is a torsor under the dilations
with no privileged origin --- the same shape as
Definition~\ref{def:resolution-torsor}.  Two consequences are worth stating,
because they are the reason this paper does not identify the two torsors.
First, the affine instance is a torsor under \(\mathbb R_{>0}\), not under
\(\mathbb Z\), so it does not yet supply the \(\langle T_\delta\rangle\cong\mathbb Z\)
structure that Theorem~\ref{thm:monodromy-difference} uses.  Second, the passage
from the affine sector to the projective one needs the trace trichotomy
\(\operatorname{tr}^2-4\det\), which Paper~0 does not develop and this paper does
not develop either.  So the instance is a calibration of the \emph{shape}, not an
identification of the objects.
\end{remark}

\subsection{Monodromy as a difference of resolutions}

\begin{theorem}[Monodromy from two through-resolutions]\label{thm:monodromy-difference}
Let $P^{\uparrow}$ and $P^{\downarrow}$ be the two resolutions of
Definition~\ref{def:resolution-torsor} obtained by passing through the critical value
along the two admissible sides.  Then the around-the-critical-value monodromy is
\[
  \bigl(P^{\downarrow}\bigr)^{-1}P^{\uparrow}=T_\delta ,
\]
the vanishing transvection.
\end{theorem}

\begin{proof}
Both resolutions identify the same two smooth fibres, and their composite is the deck
transformation of the passage around the critical value.  For a single node that deck
transformation is the Picard--Lefschetz transvection of the vanishing cycle $\delta$, by
the local theory of nondegenerate critical points.
\end{proof}

\begin{remark}
Theorem~\ref{thm:monodromy-difference} is the section's organising result.  It replaces
"insert the monodromy matrix" by "compare two resolutions", and it explains why
No-go~\ref{nogo:unipotent} had to be true: one resolution alone determines nothing,
because the torsor has no origin.
\end{remark}

\subsection{The Legendre cusp as a quadratic base change}

\begin{proposition}[Local reduction]\label{prop:legendre-reduction}
The local equation of the Legendre family at a cusp reduces, after an analytic change
of coordinates, to
\[
  UV=\lambda^{2}.
\]
The map $\lambda\mapsto s=\lambda^{2}$ is a quadratic base change of the node
$UV=s$, and consequently
\[
  M_{\mathrm{Legendre}}=T_\delta^{2}.
\]
\end{proposition}

\begin{proof}
The Legendre family has an $A_1$ node at $\lambda=0$ (and, after the standard change of
chart, at the other two cusps as well).  Pulling a node back along $\lambda\mapsto\lambda^{2}$
doubles the local monodromy, since a full turn in $s$ is covered by half a turn in
$\lambda$.  By Theorem~\ref{thm:monodromy-difference} the node monodromy is $T_\delta$,
so the pullback has monodromy $T_\delta^{2}$.
\end{proof}

\begin{remark}[What has been gained]
Proposition~\ref{prop:legendre-reduction} derives a monodromy matrix that
\S\ref{sec:classical} could only state.  The derivation uses the \emph{shape} of the
base change, not a computation with periods.  This is the sense in which the through
form is a better carrier than the operator: the operator is its shadow.
\end{remark}

\subsection{Cellular realisation on a finite annulus}

The abstract torsor has an exact finite model.

\begin{construction}[Cubical annulus and shear]\label{con:cubical-annulus}
Let $A_n$ be the cubical annulus obtained from a strip of $n$ squares by identifying the
two boundary circles with a shift, so that
\[
  \chi(A_n)=0,\qquad (b_0,b_1,b_2)(A_n)=(1,1,0).
\]
Let $D$ be the cellular shear that displaces each vertex one step along the annulus, and
let $g\colon A_n\to T_n$ be the gluing to the torus obtained by identifying the two
boundaries.
\end{construction}

\begin{proposition}[Shear induces the transvection]\label{prop:shear-transvection}
The induced map on homology is the Picard--Lefschetz transvection
\[
  D_\ast(x)=x+\langle \delta,x\rangle\,\delta ,\qquad
  M_D=\begin{pmatrix}1&1\\ 0&1\end{pmatrix},\qquad
  M_D^{T}JM_D=J ,
\]
where $J$ is the symplectic form.  Hence the shear is symplectic and realises the
primitive Dehn twist.
\end{proposition}

\begin{proof}
The shear displaces a representative of $\delta$ once around the annulus while fixing
the transverse class, which is precisely the defining formula of the transvection.  The
matrix is read off on the basis $(\delta,\gamma)$ with
$\langle\delta,\gamma\rangle=1$, and symplecticity follows from
$\langle x+\langle\delta,x\rangle\delta,\ y+\langle\delta,y\rangle\delta\rangle
=\langle x,y\rangle$ by antisymmetry and $\langle\delta,\delta\rangle=0$.
\end{proof}

\begin{theorem}[No-go: vertex closure does not imply path closure]
\label{thm:vertex-not-path}
In Construction~\ref{con:cubical-annulus} the vertex-level shear satisfies $D_0^{\,n}=I$,
while the lifted one-chain level satisfies $D_1^{\,n}\neq I$ and
$M_D^{\,n}=\begin{pmatrix}1&n\\ 0&1\end{pmatrix}\neq I$ for $n\neq0$.  Thus closure of a
finite vertex permutation is not the closure of the lifted path.
\end{theorem}

\begin{proof}
The vertex permutation is a cyclic shift of a finite set, hence has finite order dividing
$n$.  The induced homology map is unipotent with nontrivial nilpotent part, so it has
infinite order.
\end{proof}

\begin{remark}[The pair of separations]
Theorem~\ref{thm:vertex-not-path} is the finite model of
Proposition~\ref{prop:crossing-not-arrow}: a closed finite quotient can look trivial
while the lifted object is not.  \S\ref{sec:residual-catalogue} records this as one
instance of a general pattern rather than as an isolated curiosity.
\end{remark}

\begin{proposition}[Marked identity and the nodal pinch]\label{prop:marked-identity}
Let $q\colon T_n\to$ (the nodal quotient) be the pinch.  Then $q_\ast D_\ast=q_\ast$, so
the shear is invisible to the nodal reading; but the marked outgoing condition gives a
strict identity $q_+D=q_-$, which upgrades $D$ to a genuine through-resolution.
\end{proposition}

\begin{proof}
The pinch kills precisely the vanishing direction, and the shear acts on the transverse
direction by the transvection, which is trivial modulo $\delta$.  The marked statement
tracks which boundary the identity is asserted on, whence the strict identity.
\end{proof}

\subsection{Through forms as fibre products over a singular middle}

\begin{definition}[Through form]\label{def:through-form}
For a typed singular middle object $N$ and two fibres $F_\pm$ over it, the \emph{through
form} is
\[
  \langle L\mid_N R\rangle \defeq F_-\times_N F_+ .
\]
\end{definition}

\begin{proposition}[Order and compatibility are independent]\label{prop:order-compat-independent}
Conway number-admissibility of a cut $\{L\mid R\}$ and through-compatibility of
$\langle L\mid_N R\rangle$ are logically independent: either may hold while the other
fails.
\end{proposition}

\begin{proof}
Admissibility is a condition on the order structure of $L$ and $R$ alone; compatibility
is a condition on the projections to $N$.  Choosing $N$ degenerate makes compatibility
fail with admissible sides, and choosing an incompatible order on $N$ makes compatibility
hold with inadmissible sides.
\end{proof}

\begin{proposition}[Resolution fibre count]\label{prop:resolution-count}
If the critical fibre consists of $r$ regular points together with $n$ phase points,
the through form carries $r+n^{2}$ resolutions, indexed by the relative offset
$P_k\colon i\mapsto i+k$ modulo $n$.
\end{proposition}

\begin{proof}
Each regular point has one resolution; each pair of phase points may be joined, and the
join is determined by a relative offset $k$ modulo $n$, giving $n$ classes for each of
$n$ points.
\end{proof}

\begin{remark}[Duality is relational converse]
The duality appropriate to through forms is not an inverse but the relational converse
$T^{\star}=T^{\mathrm{op}}$ with contravariant composition.  This is forced by
Proposition~\ref{prop:relation-valued}: a transport that is relation-valued cannot have a
functional dual.
\end{remark}

\subsection{One feature, three readings}

\begin{theorem}[Figure-eight calibration]\label{thm:figure-eight}
There is a specific figure-eight braid word $w$ for which
\[
  -D_w=-\phi\!\left(\frac{\partial w}{\partial b}\right)
  =\det(tI-M)=t^{2}-3t+1 ,
\]
where $D_w$ is the AEG affine defect, $\phi$ is the abelianised Fox derivative
\cite{Fox1953FreeDifferential}, and $M$ the fibre monodromy.  The common value is a
feature of the ordered history and does not reconstruct the knot.
\end{theorem}

\begin{proof}[Status]
An exact computation for one word, with all three quantities evaluated independently.
The equality is a calibration on a single word, not a general identity, and the
non-reconstruction clause is essential: the polynomial is a characteristic in the sense
of Definition~\ref{def:characteristic}, with the word as its residual.
\end{proof}

\begin{proposition}[Double periods]\label{prop:double-periods}
Two visible periods $p,q$ give a phase torus $\Z/p\times\Z/q$ under the successor
$\tau(i,j)=(i+1,j+1)$, with
\[
  \#\text{orbits}=\gcd(p,q),\qquad \text{orbit length}=\lcm(p,q),\qquad
  \text{primitive winding}=\bigl(q/g,\ p/g\bigr),\quad g=\gcd(p,q).
\]
The unit-tangent lift supplies the third direction, and the orbifold meridians
\[
  m_P=(p-1,-1),\qquad m_Q=(-1,q-1),\qquad \det = pq-p-q
\]
recover the arithmetic.  Equal numerals across the three types are not identified.
\end{proposition}

\begin{proof}
An orbit of $\tau$ has length $\lcm(p,q)$ because both coordinates must return, and the
number of orbits is $pq/\lcm(p,q)=\gcd(p,q)$.  The winding is read off by dividing the
periods by $g$.  The determinant computation is direct.
\end{proof}

\subsection{Return is not period}

\begin{nogo}[No intrinsic periodic dynamics]\label{nogo:return-not-period}
There is a finite diagram containing the source-free components
$\mathrm{constant};\mathrm{discard}$ and
$\mathrm{constant};\mathrm{neg};\mathrm{discard}$ whose least boundary-return lengths are
$2$ and $3$ respectively; but monotonicity of the completed-past rank forces the verdict
\emph{NotRepresentable} for intrinsic periodic dynamics.  Hence a boundary-return word is
not a periodic orbit.
\end{nogo}

\begin{proof}
The completed past of a state grows monotonically along a checked execution, so no state
can recur.  The return lengths are therefore properties of words, not periods of a
dynamical system; a genuine period requires a feedback operator that the syntax does not
supply.
\end{proof}

\begin{remark}
No-go~\ref{nogo:return-not-period} is the most frequently violated boundary in this
material.  It is why Part~II treats iteration as a syntactic structure
(\S\ref{sec:boundaries}) rather than as a semantic one.
\end{remark}

\subsection{Summary}

\begin{enumerate}
\item A crossing is a cospan with a resolution torsor
  (Definition~\ref{def:through-crossing}, Proposal~\ref{prop:crossing-not-arrow},
  Definition~\ref{def:resolution-torsor}).
\item Monodromy is the difference of two resolutions, and the Legendre monodromy follows
  from the quadratic base change (Theorem~\ref{thm:monodromy-difference},
  Proposition~\ref{prop:legendre-reduction}).
\item A primitive transvection is realised by an exact cellular shear, and vertex
  closure does not imply path closure
  (Proposition~\ref{prop:shear-transvection}, Theorem~\ref{thm:vertex-not-path}).
\item Through forms are fibre products over a singular middle, whose duality is
  relational converse (Definition~\ref{def:through-form},
  Propositions~\ref{prop:order-compat-independent}, \ref{prop:resolution-count}).
\item One feature can be read three ways on one word, and a boundary-return word is still
  not a period (Theorem~\ref{thm:figure-eight},
  Proposition~\ref{prop:double-periods}, No-go~\ref{nogo:return-not-period}).
\end{enumerate}
